\chapter{Supply Trades}\label{ch:s2:supply-trades}

Treasuries cheapen in the days before a large auction and recover after, although every auction's date and size are announced in advance. Lou, Yan
and Zhang documented the pattern and put the Treasury's hidden cost at 9 to 18 basis points of each auction's size. Dealers must absorb the supply,
and they are paid for it. From 2010 to September 2026, the ten-year yield rose 1.98 basis points on average over the five days before its 301 auctions
($t$ = 3.0), and the thirty-year 2.56 before its 234 ($t$ = 3.6). On this chapter's synthetic market, a trade that is short duration before each
auction and long after earns 3.44 basis points of yield per auction ($t$ = 3.1 over 239), a Sharpe ratio of 0.70 a year, while any single auction's
result swings by tens of basis points. The build is \texttt{firm.supplytrade}.

\section{Auction concessions}

\begin{definition}[Auction concession]\label{def:s2:supply-trades:concession}
An \emph{auction concession}\index{auction concession} is the cheapening of a security, and of those close to it, in the days before a new issue is
sold: the price the market charges for absorbing a known supply, recovered in part once the issue has been distributed.
\end{definition}

The new-issue concession of a single corporate bond (Book~2, chapter~21) is the same idea for one issuer. Fleming, Nguyen and Rosenberg studied US
Treasury dealer positions from 1990 to 2020. Issuance was the main driver of dealers' weekly inventory changes. Those changes were only partly offset
in adjacent weeks and not significantly hedged with futures, and dealers were paid for the inventory risk by later price appreciation. After the
crisis, as balance sheet grew dearer, dealers took smaller positions and laid them off faster. Investment funds' growing participation in auctions
has reduced the compensation.

The pattern is visible in the public data (\cref{fig:s2:supply-trades:path}). Group each note and bond auction from 2010 to 2026 with its nearest
benchmark and measure that benchmark's constant-maturity yield around the auction day. Pooled over 1,338 auctions, the yield is 1.47 basis points
higher on auction day than ten days before, and 5 days later it has given back most of that: 0.40 above. The pattern differs by maturity:

\begin{center}
\small
\begin{tabular}{@{}lccccc@{}}
\toprule
2010--2026 & auctions & 5 days before (bp) & $t$ & 5 days after (bp) & $t$\\
\midrule
2-year & 354 & $-0.40$ & $-0.9$ & $-0.62$ & $-1.3$\\
5-year & 249 & 1.03 & 1.5 & $-2.16$ & $-3.1$\\
7-year & 200 & $-0.53$ & $-0.7$ & $-2.02$ & $-2.1$\\
10-year & 301 & 1.98 & 3.0 & $-0.28$ & $-0.4$\\
30-year & 234 & 2.56 & 3.6 & $-0.81$ & $-1.1$\\
\bottomrule
\end{tabular}
\end{center}

The long end cheapens before, and the 5- and 7-year richen after. The 2-year shows nothing. Sorted by offering size, the ten-year's rise before the
auction grows with size: 1.48 basis points for auctions averaging \$14.5 billion, 1.69 for \$22.3 billion, 2.95 for \$36.2 billion. Its
fall after does not: $-1.64$, $-0.24$ and $+1.44$. The constant-maturity yield is an end-of-day series, so on auction day it already includes the
auction's result.

\begin{omfigure}
\begin{tikzpicture}
\begin{axis}[width=10cm,height=5.4cm,xlabel={\small trading days from the auction},ylabel={\small yield change (bp)},tick label style={font=\footnotesize},
  xmin=-10,xmax=10,ymin=-0.5,ymax=2,grid=major,grid style={gray!25},table/col sep=comma]
\addplot[omThm,thick,mark=*,mark size=1.2pt] table[x=day,y=bp]{figdata/strategies-2/13-supply-trades/path.csv};
\addplot[black,thin] coordinates {(0,-0.5) (0,2)};
\end{axis}
\end{tikzpicture}

\omcaption{The average path of the benchmark Treasury yield around 1,338 note and bond auctions, 2010--2026, relative to ten trading days before,
from TreasuryDirect's auction records and the Board's constant-maturity yields on FRED. Derived statistics only. Data:
\texttt{s2\_fetch\_auctions}.}\label{fig:s2:supply-trades:path}
\end{omfigure}

\section{New-issue and on-the-run premia}

\begin{definition}[On-the-run premium]\label{def:s2:supply-trades:otr}
The \emph{on-the-run premium}\index{on-the-run premium} is the amount by which the most recently issued security at a maturity trades rich to older
securities of similar maturity, for its liquidity and its value in repo; it shrinks as the issue ages and a newer one takes its place.
\end{definition}

The premium shapes what a \omterm{def:s2:supply-trades:trade}{supply trade} holds. The new issue is cheap to the curve before its auction and becomes the rich on-the-run after it. The
issue it replaces loses its status and cheapens. The Board's fitted curve (chapter~10) leaves out the on-the-run and first off-the-run issues for this
reason. A roll trade sells the old on-the-run issue and buys the new one around the auction. It bets that the premium moves from one to the other
more slowly than the market prices.

\section{Corporate supply}

Corporate issuance works the same way with smaller, less predictable deals. Issuers announce in the morning and price by the afternoon, and dealers
and investors buy at a concession to the issuer's existing bonds. Heavy issuance weeks cheapen the whole sector. Before a large deal the issuer's own
bonds and close substitutes are sold, or hedged with swaps and Treasuries, and the hedges come off afterwards. The effect is predictable when the
calendar is. It is also crowded: many desks know the calendar.

\begin{definition}[Supply trade]\label{def:s2:supply-trades:trade}
A \emph{supply trade}\index{supply trade} takes the opposite side of known, scheduled supply: short (or underweight) the securities about to be issued
and their close substitutes before the sale, long after it, paid by the issuer's and the dealers' need to distribute the supply.
\end{definition}

\section{Trading around the calendar}

\texttt{firm.supplytrade} makes the concession explicit (\cref{lst:s2:supply-trades:model}). A benchmark yield carries 5.5 basis points a day of
noise and a monthly auction of \$20 to \$45 billion. Dealers' capacity to absorb each auction varies (a lognormal with dispersion 0.4). In the five
days before an auction the yield rises by 0.07 basis points per billion at average capacity, 2.60 on average, and in the five days after it gives back
70\% of the rise. The trade is short duration from five days before to the auction's close and long for five days after, paying 0.1 basis points of
yield per leg in futures.

Over twenty years and 239 auctions it earns 3.44 basis points per auction ($t$ = 3.1): 1.71 from the leg before and 1.72 from the leg after. At 12
auctions a year that is a Sharpe ratio of 0.70, and 0.50 for either leg alone. Sorted by size, the average rises from 1.66 basis points in the
smallest third to 4.94 in the middle and 3.72 in the largest. Each billion of size adds 0.24 basis points to the trade on average
(\cref{fig:s2:supply-trades:scatter}). The noise is far larger than the effect: single auctions range from $-47.9$ to $+52.7$ basis points. Dealers'
capacity, which the trader does not see, decides as much as size. The slope on size over capacity is 0.16 per unit.

\begin{omfigure}
\begin{tikzpicture}
\begin{axis}[width=10cm,height=5.8cm,xlabel={\small auction size (\$ billion)},ylabel={\small trade P\&L (bp of yield)},tick label style={font=\footnotesize},
  xmin=19,xmax=46,ymin=-55,ymax=60,grid=major,grid style={gray!25},table/col sep=comma]
\addplot[only marks,mark=*,mark size=1pt,omThm,opacity=0.6] table[x=size,y=pnl]{figdata/strategies-2/13-supply-trades/scatter.csv};
\addplot[black,thick] table[x=size,y=fit]{figdata/strategies-2/13-supply-trades/fit.csv};
\end{axis}
\end{tikzpicture}

\omcaption{The synthetic \omterm{def:s2:supply-trades:trade}{supply trade}'s P\&L per auction against the auction's size, 239 auctions over twenty years, with the least-squares line (0.24
basis points per billion). Data: \texttt{s2\_supply.results}.}\label{fig:s2:supply-trades:scatter}
\end{omfigure}

\section{Strategy files}

\begin{strategyfile}{Pre-auction short, post-auction long}\label{strat:s2:supply-trades:auction}
\sfield{Who pays you, and why} The Treasury, through the concession it pays to distribute supply; dealers with limited capacity.
\sfield{Instruments and venues} Treasury futures or the benchmark and its neighbours; the when-issued market.
\sfield{Signal} The auction calendar and size; dealer positioning where visible.
\sfield{Sizing and execution} Small per auction, many auctions; futures for low cost.
\sfield{Costs} Futures fees and bid--ask; roll.
\sfield{How it dies} More non-dealer demand at auctions; crowding by others trading the calendar.
\sfield{Horizon, capacity, infrastructure} Days; the calendar.
\sfield{Backtest honestly} Yields at the actual trade times, not end of day.
\sfield{Sources} Lou, Yan and Zhang (2013); 2010--2026 yields: $+1.98$ basis points before 10-year auctions; this chapter: 3.44 per auction, Sharpe
ratio 0.70.
\end{strategyfile}

\begin{strategyfile}{On-the-run roll}\label{strat:s2:supply-trades:roll}
\sfield{Who pays you, and why} Holders who pay for on-the-run liquidity and repo value.
\sfield{Instruments and venues} The new and the old on-the-run issues.
\sfield{Signal} The yield spread between them against its pattern over the cycle.
\sfield{Sizing and execution} DV01-matched; the short financed in reverse repo.
\sfield{Costs} Specialness of the issue shorted.
\sfield{How it dies} A flight to liquidity widens the premium.
\sfield{Horizon, capacity, infrastructure} An auction cycle.
\sfield{Backtest honestly} Special repo rates.
\sfield{Sources} No performance figure verified.
\end{strategyfile}

\begin{strategyfile}{Corporate new-issue concession}\label{strat:s2:supply-trades:corporate}
\sfield{Who pays you, and why} Issuers who price new bonds below their existing ones to sell them quickly.
\sfield{Instruments and venues} New corporate bonds at allocation; the issuer's secondary bonds.
\sfield{Signal} The concession against the issuer's curve.
\sfield{Sizing and execution} Allocations; hedged with Treasuries or swaps.
\sfield{Costs} Allocations smaller than orders; secondary spreads.
\sfield{How it dies} Crowded books that shrink the concession.
\sfield{Horizon, capacity, infrastructure} Days; relationships with underwriters.
\sfield{Backtest honestly} Actual allocations, not orders.
\sfield{Sources} No performance figure verified.
\end{strategyfile}

\begin{strategyfile}{Month-end index extension}\label{strat:s2:supply-trades:extension}
\sfield{Who pays you, and why} Index-tracking investors who must add duration when the index's duration rises at month-end, as new issues enter it.
\sfield{Instruments and venues} Long bonds or futures before month-end.
\sfield{Signal} The index's announced extension.
\sfield{Sizing and execution} Positions before the rebalancing, unwound after.
\sfield{Costs} Futures fees.
\sfield{How it dies} Many desks trade the same calendar.
\sfield{Horizon, capacity, infrastructure} Days.
\sfield{Backtest honestly} Extension figures as published before month-end.
\sfield{Sources} No figure verified.
\end{strategyfile}

\section{Tutorial: paid to absorb}\label{tut:s2:supply-trades:tut}

\begin{tutorial}
\textbf{Goal.} Simulate an auction calendar with planted concessions, trade it, and measure the real pattern around 2010--2026 auctions.
\textbf{End state:} the tables and the two figures.
\begin{tutsteps}
\item \textbf{Auctions and the trade}.
\omcode{code/firm/supplytrade/firm_supplytrade.py}{38}{62}{A calendar with concessions over size and capacity, and the trade's legs.}\label{lst:s2:supply-trades:model}
\item \textbf{Results}.
\omcode{code/strategies-2/13-supply-trades/python/s2_supply.py}{35}{51}{By leg, by size, and against size.}\label{lst:s2:supply-trades:results}
\item \textbf{Run} \texttt{s2\_fetch\_auctions.py} once, then \texttt{results()} and \texttt{fig\_supply.py}.
\end{tutsteps}
\textbf{What to change next.} Give the trader a noisy signal of dealer capacity and trade only when it is low; add an \omterm{def:s2:supply-trades:otr}{on-the-run premium} that moves at
each auction; use intraday yields at the auction time.
\end{tutorial}

\section{Build: supply trades}\label{bld:s2:supply-trades:supplytrade}

\begin{build}
\bfield{Purpose} An auction calendar with concessions proportional to size over dealer capacity, and the trade around it.
\bfield{Interface} \texttt{SupplyConfig(\dots)}, \texttt{simulate\_auctions(cfg)}, \texttt{supply\_trade(sim, cfg, pre, post)}.
\bfield{Rules} The concession builds over the five days before and is partly given back over the five after; costs per leg.
\bfield{Acceptance tests} \texttt{code/firm/supplytrade/tests/}: without noise the trade earns the concession and its giveback; costs per leg;
concessions scale with size over capacity.
\bfield{Stretch} Several maturities; reopenings; an \omterm{def:s2:supply-trades:otr}{on-the-run premium}.
\end{build}

\begin{omsources}
\item D.~Lou, H.~Yan and J.~Zhang, ``Anticipated and repeated shocks in liquid markets'', \emph{Review of Financial Studies} 26(8), 2013.
\item M.~Fleming, G.~Nguyen and J.~Rosenberg, ``How do Treasury dealers manage their positions?'', Federal Reserve Bank of New York Staff Report 299,
revised 2024.
\item TreasuryDirect auction records; Board of Governors H.15 yields, via FRED.
\end{omsources}

\section{Exercises}

\begin{exercise}[$\star$]\label{exo:s2:supply-trades:1}
An auction of \$40 billion at average capacity carries a concession of 0.07 basis points per billion. How large is it?
\end{exercise}

\begin{exercise}[$\star$]\label{exo:s2:supply-trades:2}
A trade earns 3.44 basis points per auction with a standard deviation of about 17 basis points, 12 times a year. What is its Sharpe ratio a year?
\end{exercise}

\begin{exercise}[$\star$]\label{exo:s2:supply-trades:3}
Lou, Yan and Zhang put the hidden cost at 9 to 18 basis points of the auction size. What is that on a \$40 billion auction?
\end{exercise}

\begin{exercise}[$\star\star$]\label{exo:s2:supply-trades:4}
Why can a known, announced supply move prices at all?
\end{exercise}

\begin{exercise}[$\star\star$]\label{exo:s2:supply-trades:5}
Why might the compensation for absorbing auctions have fallen since 2008?
\end{exercise}

\begin{exercise}[$\star\star$]\label{exo:s2:supply-trades:6}
Why does an end-of-day yield understate the concession on auction day?
\end{exercise}

\begin{exercise}[$\star\star\star$]\label{exo:s2:supply-trades:7}
\emph{Coding.} Rerun with \texttt{SupplyConfig(giveback=0.0)}. What happens to each leg, and what does it say about which leg is the dealers'
compensation?
\end{exercise}

\begin{exercise}[$\star\star\star$]\label{exo:s2:supply-trades:8}
\emph{Find the flaw.} ``The last three auctions each lost us 20 basis points, so the supply effect has gone.''
\end{exercise}

\section{Problem: Paid to Absorb}

\begin{problem}[{Weekend problem --- trading the auction calendar}]\label{pb:s2:supply-trades:1}
The chapter's synthetic calendar, 2010--2026 auctions and the public record.

\medskip\noindent\textbf{Part I --- Concessions.}
\begin{enumerate}
\item Define an \omterm{def:s2:supply-trades:concession}{auction concession} and the \omterm{def:s2:supply-trades:otr}{on-the-run premium}.
\item What did Lou, Yan and Zhang find, and what did they estimate?
\item What did Fleming, Nguyen and Rosenberg find about dealers?
\item Why does the Board's curve exclude the newest issues?
\end{enumerate}

\medskip\noindent\textbf{Part II --- The real pattern.}
\begin{enumerate}[resume]
\item Give the average path around auctions.
\item Which maturities cheapen before, and which richen after?
\item How does the 10-year's pattern vary with size?
\item What does the end-of-day series miss?
\end{enumerate}

\medskip\noindent\textbf{Part III --- The synthetic trade.}
\begin{enumerate}[resume]
\item Define the \omterm{def:s2:supply-trades:trade}{supply trade}.
\item Describe the synthetic concession.
\item Give the trade's mean per auction, $t$-statistic and Sharpe ratio, and each leg's.
\item Why is any one auction's result so noisy?
\end{enumerate}

\medskip\noindent\textbf{Part IV --- The verdict.}
\begin{enumerate}[resume]
\item State the \emph{named result}: the auction trade's return per auction and its dependence on auction size.
\item What does dealers' capacity add?
\item How would you trade corporate supply?
\item What is the month-end extension trade?
\item How would you backtest the auction trade honestly?
\item Which strategy file depends on repo?
\item How does this chapter relate to chapter~10?
\item In one sentence: who pays the supply trader?
\end{enumerate}
\end{problem}

\section{Interview questions}

\begin{interviewq}[$\star$ \iqroles{trader}]\label{iq:s2:supply-trades:1}
Why would Treasuries cheapen before an auction everyone knows about?
\end{interviewq}

\begin{interviewq}[$\star\star$ \iqroles{researcher}]\label{iq:s2:supply-trades:2}
How would you measure the \omterm{def:s2:supply-trades:concession}{auction concession}, and what would you control for?
\end{interviewq}

\begin{interviewq}[$\star\star$ \iqroles{trader}]\label{iq:s2:supply-trades:3}
You are a dealer bidding in tomorrow's ten-year auction. How do you hedge?
\end{interviewq}

\begin{interviewq}[$\star\star$ \iqroles{risk}]\label{iq:s2:supply-trades:4}
A desk runs the auction trade each month. What is its risk?
\end{interviewq}

\begin{interviewq}[$\star\star$ \iqroles{developer}]\label{iq:s2:supply-trades:5}
Design a data pipeline for auction announcements and results.
\end{interviewq}

\begin{interviewq}[$\star\star\star$ \iqroles{researcher}]\label{iq:s2:supply-trades:6}
With daily yield noise of $\sigma$ basis points and a concession $c$ built over $k$ days, how many auctions does it take to detect $c$ with a $t$ of 2?
\end{interviewq}
